\subsection{12.40: unbounded Brauer-degree parts at fixed Sylow order}
\label{cand:12.40}
\begin{proposition}
For every \(k\ge3\), some odd positive integer \(a\) gives an
irreducible \(2\)-Brauer character \(\varphi\) of
\(G=\mathrm{PSL}_6(5^a)\) with
\[
 |G|_2=8192,\qquad \varphi(1)_2=2^k.
\]
Thus no finite-valued bound \(\varphi(1)_p\le f(|G|_p)\) exists,
already for the fixed prime \(p=2\).
\end{proposition}
\begin{proof}
Bray's projective-point module \cite[Main Theorem]{bray-cohomology}
has an absolutely irreducible heart in characteristic two of dimension
\[
 F(q)=\frac{q^6-1}{q-1}-2=q^5+q^4+q^3+q^2+q-1.
\]
Here \(q=5^a\). To recall why the assertion applies, let \(P\) be the
projective permutation module over an algebraically closed field of
characteristic two, \(U\) its constant line, and \(W\) its augmentation
kernel. Its even permutation degree puts \(U\subseteq W\).
Perfectness of \(\mathrm{SL}_6(q)\) excludes a trivial submodule of
\(W/U\): its preimage would be an extension of two trivial modules,
which splits, contradicting the one-dimensional invariant space of \(P\).
The permutation form makes \(W/U\) self-dual, so there is no trivial
quotient either. Each nontrivial simple factor has dimension at least
\(q^5-1\): on the elementary abelian unipotent subgroup \(\F_q^5\)
it has a nontrivial weight, and the Levi subgroup
\(\mathrm{SL}_5(q)\) permutes all nontrivial weights transitively.
The unipotent subgroup normally generates \(\mathrm{SL}_6(q)\), so
it cannot act trivially on a nontrivial simple factor.
Since \(2(q^5-1)>F(q)\), there is only one nontrivial factor.
Absence of trivial submodules and quotients now makes the heart simple.
Scalars fix projective points, so it descends to \(\mathrm{PSL}_6(q)\).

For odd \(a\), \(q\equiv5\pmod8\) and \(q\equiv-1\pmod3\);
hence \(\gcd(6,q-1)=2\). The order formula is
\[
 |G|=\frac{q^{15}}2\prod_{i=2}^{6}(q^i-1).
\]
The five factors have \(2\)-adic valuations \(3,2,4,2,3\).
Their sum, less the central factor of two, is \(13\), independently
of \(a\).

It remains to obtain arbitrarily prescribed degree valuations at
actual prime powers. Repeated squaring gives
\(v_2(5^{2^j}-1)=j+2\), so for \(k\ge3\),
\[
 5^{2^{k-2}}\equiv1+2^k\pmod{2^{k+1}}.
\]
For odd \(x\), \(F'(x)\) is odd, and the binomial theorem gives
\[
 F(x+2^k)\equiv F(x)+2^k\pmod{2^{k+1}}.
\]
Therefore, whenever odd \(a\) satisfies \(2^k\mid F(5^a)\),
exactly one of \(a\) and \(a+2^{k-2}\) gives divisibility by
\(2^{k+1}\); the other gives valuation exactly \(k\).
Start with \(a=1\), since \(F(5)=3904\) is divisible by eight,
and continue using the divisible choice. The unused choice at each
step supplies the required odd exponent and exact valuation.
\end{proof}

The representation family and its irreducibility are prior mathematics;
the fixed group-order valuation and unbounded degree valuation are the
deduction proposed here. No assertion about ordinary character degrees
is being made. See \artifact{research/12.40-proof.md};
finite module and integer controls corroborate the formulas without
replacing the arbitrary-\(k\) argument.
